\section{Cross-Cutting Findings}\label{sec:findings}

\subsection{Latency: how long silence lasts}\label{sec:latency}

\begin{figure}[t]
\centering
\begin{tikzpicture}[xscale=0.078]
  % log-ish horizontal bars: silence span in days (not to scale below 1d)
  \foreach \y/\len/\name/\mech in {
      0/60/SSD backup (TCC sandbox)/architecture-level assumption + compromised forensics,
      -1/7/watchdog self-death/monitoring layer's own silence,
      -2/5.5/observer path; backup anti-pattern/deployment topology; copy-pasted suppression,
      -3/0.55/gateway death (9h); reserved-file mute (13h)/alert-chain dependency; runtime semantics,
      -4/1.6/digest 502; cron omission/error dilution; operational omission}{
    \draw[fill=red!12, draw=red!60!black, rounded corners=1pt]
      (0,\y*0.62) rectangle (\len,\y*0.62+0.38);
    \node[anchor=west, font=\scriptsize] at (\len+1.2,\y*0.62+0.19)
      {\name\ \ \textcolor{gray}{---\ \mech}};
  }
  \draw[-{Stealth}, thin] (0,-3.2) -- (64,-3.2);
  \foreach \x in {0,10,...,60} \node[font=\tiny, below] at (\x,-3.2) {\x};
  \node[font=\scriptsize, below=3.5mm] at (30,-3.2) {silence span (days)};
\end{tikzpicture}
\caption{Silence span by incident group. Latency tracks the mechanism
layer, not code complexity: the long tail lives where no test runs.}
\label{fig:latency}
\end{figure}

Latency tracks the \textbf{mechanism layer, not code complexity}
(Fig.~\ref{fig:latency}). Code-level bugs die young: unit tests and the
next run catch them. What survives for weeks lives where no test
runs---in deployment topology, OS policy, monitoring-of-monitoring, and
the gap between declared and runtime state. Latency is therefore a
\emph{measure of observational distance}: the further a mechanism sits
from any existing observer, the longer it lives. We suggest
silence-latency percentiles as a reportable reliability metric for agent
systems, complementary to MTTR---our experience is that for silent
failures, time-to-\emph{detect} dominates time-to-repair by one to two
orders of magnitude.

\subsection{Discovery: who finally notices}\label{sec:discovery}

\begin{table}[t]
\centering\small
\caption{Discovery channels across the corpus (qualitative shares).}
\label{tab:discovery}
\begin{tabular}{@{}p{52mm}p{18mm}p{52mm}@{}}
\toprule
\textbf{Channel} & \textbf{Share} & \textbf{Notes} \\
\midrule
\textbf{Human user-view} (reading actual pushed output; weekly
observation ritual) & \textbf{$\sim$70\%} & ``this digest looks
shallow'', ``why two windows?'', ``didn't receive yesterday's
analysis'' \\
\addlinespace[2pt]
Target-environment execution (dev green, prod reveals) & high & all
\classref{A}; several \classref{B} \\
\addlinespace[2pt]
Self-observation (governance auditing governance; observer critiquing
output) & rising & added mid-study; caught its own executor bug \\
\addlinespace[2pt]
Unit tests / preflight & $\approx$0 for this corpus & by selection:
anything they caught never became a silent incident \\
\bottomrule
\end{tabular}
\end{table}

The last row of Table~\ref{tab:discovery} is partly tautological---our
corpus selects for what tests missed---but the magnitude of the first row
is not. The system runs 4{,}286 tests, 827 governance checks, and a
19-point preflight, all green through most of these incidents; the single
most productive detector of silent failure was \textbf{a human looking at
the product as a user}. We institutionalized this as a weekly 30-minute
observation ritual (no coding allowed; four dimensions: alert noise, push
latency, information density, response quality)---and it continued to
out-detect the automated stack. For practitioners, the implication is
blunt: \emph{user-view observation is a first-class observability signal
and deserves calendar time}; for researchers, the open problem is
mechanizing even part of what the human eye does here. The LLM-as-judge
literature gives grounds for optimism---strong LLM judges reach
${>}80\%$ agreement with human preference on open-ended response
quality~\cite{zheng2023judge}---and our own step in that direction, a
daily LLM ``observer'' that critiques yesterday's outputs against quality
heuristics, found real regressions (including a fabrication, and
including two bugs in \emph{itself}). The qualifier our experience adds
to that literature: an LLM judging a \emph{system's output for silent
failure} is itself an LLM component of that system, inheriting every
class in this taxonomy---ours shipped with a \classref{B} path bug and a
sampling artifact that made it hallucinate a truncation---so the judge
needs the same governance, provenance hygiene, and sabotage validation as
the components it judges.

\subsection{The trigger--amplifier--concealer structure}\label{sec:tac}

Nearly every postmortem decomposed into three causally distinct layers:

\begin{center}\small
\begin{tabular}{@{}lll@{}}
\toprule
\textbf{Layer} & \textbf{Role} & \textbf{Examples} \\
\midrule
trigger & the external spark & a surrogate byte; an omitted LLM output
line; a transient EPERM \\
amplifier & the architectural flaw that spreads it & stdout logging into
command substitution; \\
& & positional parsing; 18 copy-pasted suppression idioms \\
concealer & the absence that hides it & status file lying ``ok'';
fail-open guard; quiet-hours filter; \\
& & forensic tool silently denied \\
\bottomrule
\end{tabular}
\end{center}

The practical force of the decomposition is prescriptive: \textbf{a fix
that addresses only the trigger is cosmetic}. Triggers are unbounded (the
environment will always produce another malformed byte, another quirk);
amplifiers and concealers are finite and owned by the architecture. Every
recurrence in our corpus (\S\ref{sec:bugclass}) traces to a fix that
stopped at the trigger. Conversely, the highest-leverage single fixes
were amplifier-level (one \texttt{>\&2}; one shared helper replacing 20
idiom copies) and concealer-level (fail-loud with cause; forensic stderr
tagging).

\subsection{Silent failure is a bug class, not a bug}\label{sec:bugclass}

The MR-4 manifestation counter (\S\ref{sec:counter}) kept advancing at a
roughly constant rate across the study---but never twice in the same
form. Early manifestations were classic swallowing; middle-period ones
were dilution and fabrication; late ones included \emph{a fix introducing
a deeper silence} (a retry helper whose exit-code propagation silently
killed twenty \texttt{set~-e} callers mid-script---the fix for a silent
failure became a worse one), \emph{a correct fix to the wrong bug} (an
ownership repair that was real but irrelevant to the EPERM it targeted,
exposed only because the metric refused to move), and \emph{vacuous
verification} (the 67 checks executing empty strings). The lesson we draw
is structural: silent failure cannot be enumerated and fixed; it can only
be \textbf{governed}---by meta-rules that constrain whole mechanism
families, scanners that enforce the meta-rules mechanically, and
validation that the guards themselves are alive (\S\ref{sec:defense}).

\subsection{Defense maturation: point fix $\rightarrow$ meta-rule
$\rightarrow$ scanner}\label{sec:maturation}

\begin{figure}[t]
\centering
\begin{tikzpicture}[node distance=3mm and 8mm]
  \node[stage, minimum width=34mm, minimum height=15mm] (fix)
    {\textbf{1. point fix}\\\scriptsize repair this bug};
  \node[amp, minimum width=34mm, minimum height=15mm, right=of fix] (rule)
    {\textbf{2. meta-rule}\\\scriptsize named, cross-case rule\\\scriptsize (23 by study end)};
  \node[good, minimum width=34mm, minimum height=15mm, right=of rule] (scan)
    {\textbf{3. scanner}\\\scriptsize mechanized check in\\\scriptsize CI + daily audit (15)};
  \draw[flow] (fix) -- (rule);
  \draw[flow] (rule) -- (scan);
  \node[lbl, below=2mm of fix]  {recurs: lesson exists\\only as a diff\\\textcolor{red!60!black}{(8 sites, 2 days later)}};
  \node[lbl, below=2mm of rule] {memory-bound:\\constrains authors\\who remember it};
  \node[lbl, below=2mm of scan] {recurrence structurally\\impossible\\\textcolor{green!40!black}{(0 recurrences to date)}};
\end{tikzpicture}
\caption{Defense maturation path. Every meta-rule that reached step~3 has
zero recorded recurrences; every recurrence in the corpus involves a
lesson that stopped at step~1.}
\label{fig:maturation}
\end{figure}

Defense effectiveness in our record correlates with how far a lesson
traveled along a three-step path (Fig.~\ref{fig:maturation}):

\begin{enumerate}
\item \textbf{Point fix} --- repair this bug. \emph{Empirically
insufficient:} an import-omission bug fixed at one site recurred at eight
sites via an automated injector \textbf{two days later}, because the
lesson existed only as a diff.
\item \textbf{Meta-rule} --- write the lesson as a named, cross-case rule
(23 such rules by study end; e.g., diagnostics-to-stderr; key-based
parsing; alert-path independence; declared-state convergence;
reserved-file unwritability). Necessary but memory-bound: rules constrain
authors who remember them.
\item \textbf{Mechanized scanner} --- encode the meta-rule as a
repository-wide check that runs in CI and the daily audit (15 such
scanners by study end: cross-OS quirks, path consistency, heredoc import
closure, monitor self-alarm compliance, cross-environment path
resolution, \dots). Only at this step does recurrence become structurally
impossible rather than culturally discouraged.
\end{enumerate}

Every meta-rule that reached step~3 has zero recorded recurrences as of
the study cutoff; every recurrence in the corpus involves a lesson that
had stopped at step~1.

\subsection{Audit is a regression engine, not a prediction
engine}\label{sec:audit}

Mid-study we audited our own defense framework against the first 15
incidents with three questions per incident: Q1---could the audit, as it
existed \emph{before} the incident, have caught it? Q2---why not?
Q3---do the guards added \emph{after} it block the same class going
forward? Results:

\begin{center}\small
\begin{tabular}{@{}lr@{}}
\toprule
\textbf{Metric} & \textbf{Value} \\
\midrule
Ex-ante prevention (Q1 fully) & \textbf{0 / 15 = 0\%} \\
Partial early warning & 2 / 15 = 13\% \\
Ex-post regression blocking (Q3 $\geq$ half) & \textbf{13 / 15 = 87\%} \\
Root cause of misses: a dimension the audit had never conceived &
12 / 15 = 80\% \\
\bottomrule
\end{tabular}
\end{center}

The 0\% is not an indictment of the audit---it is its \emph{job
description}. 80\% of misses were ``blank categories'': dimensions no
invariant had ever contemplated, which no amount of diligence within the
existing dimension set would have covered. Audits, like regression test
suites, encode the past. The honest engineering posture is to
(a)~maximize the regression rate (ours: 87\%, with sabotage validation
pushing reliability of that number), (b)~accept that prevention of
\emph{novel} mechanism classes comes from elsewhere---user-view
observation, adversarial review (``what could break that we would not
notice?''), and target-environment exposure---and (c)~measure the
\emph{conversion latency} from novel incident to mechanized guard, which
the three-step path (\S\ref{sec:maturation}) is designed to minimize. A
complementary adversarial audit (16 destruction scenarios injected
against the live repository: 10 replaying known incidents, 6 probing
suspected blind spots) scored 16/16 after the blind-spot batch drove its
own round of guard additions---a useful exercise precisely because its
first run did not score 16/16.
